\documentclass[12pt, a4paper]{article}
\usepackage{../../../myconfig} % Gọi file cấu hình từ ngoài gốc vào

% ==========================================
% BẮT ĐẦU NỘI DUNG TÀI LIỆU
% ==========================================
\begin{document}

% Truyền 3 tham số: {Nhãn cam} {Chữ to chính giữa} {Chữ ở góc phải Header}
\titlebox{Chủ đề: Hàm số}{Lượng giác}{Hàm số: Lượng giác}

% --- CÂU 1 ---
\questionLinesManual{0}{1}{Tập nghiệm của phương trình \( \sin x = 1 \) là}{
    \begin{tasks}(2)
        \task \( S = \left\{ \dfrac{\pi}{2} + k\pi \mid k \in \mathbb{Z} \right\} \)
        \task \( S = \left\{ \dfrac{\pi}{2} + k2\pi \mid k \in \mathbb{Z} \right\} \)
        \task \( S = \{ k\pi \mid k \in \mathbb{Z} \} \)
        \task \( S = \{ k2\pi \mid k \in \mathbb{Z} \} \)
    \end{tasks}
}

% --- CÂU 2 ---
\questionLinesManual{0}{2}{Tập nghiệm của phương trình \( \sin x = 0 \) là}{
    \begin{tasks}(2)
        \task \( S = \left\{ \dfrac{\pi}{2} + k\pi \mid k \in \mathbb{Z} \right\} \)
        \task \( S = \{ k\pi \mid k \in \mathbb{Z} \} \)
        \task \( S = \left\{ \dfrac{\pi}{2} + k2\pi \mid k \in \mathbb{Z} \right\} \)
        \task \( S = \{ k2\pi \mid k \in \mathbb{Z} \} \)
    \end{tasks}
}

% --- CÂU 3 ---
\questionLinesManual{0}{3}{Cho \( \sin \alpha = \dfrac{3}{5} \) và \( \dfrac{\pi}{2} < \alpha < \pi \). Giá trị của \( \cos \alpha \) là:}{
    \begin{tasks}(4)
        \task \( -\dfrac{4}{5} \)
        \task \( -\dfrac{4}{5} \)
        \task \( +\dfrac{4}{5} \)
        \task \( -\dfrac{16}{25} \)
    \end{tasks}
}

% --- CÂU 4 ---
\questionLinesManual{0}{4}{Cho \( \sin x = \dfrac{2}{3} \). Giá trị của biểu thức \( A = \cos 2x + 1 \) bằng}{
    \begin{tasks}(4)
        \task \( \dfrac{10}{9} \)
        \task \( \dfrac{26}{9} \)
        \task \( \dfrac{10}{3} \)
        \task \( -\dfrac{2}{3} \)
    \end{tasks}
}

% --- CÂU 5 ---
\questionLinesManual{0}{5}{Tìm tập xác định \( D \) của hàm số \( y = \tan 2x \):}{
    \begin{tasks}(2)
        \task \( D = \mathbb{R} \setminus \left\{ \dfrac{\pi}{4} + k2\pi \mid k \in \mathbb{Z} \right\} \)
        \task \( D = \mathbb{R} \setminus \left\{ \dfrac{\pi}{2} + k\pi \mid k \in \mathbb{Z} \right\} \)
        \task \( D = \mathbb{R} \setminus \left\{ \dfrac{\pi}{4} + k\pi \mid k \in \mathbb{Z} \right\} \)
        \task \( D = \mathbb{R} \setminus \left\{ \dfrac{\pi}{4} + \dfrac{k\pi}{2} \mid k \in \mathbb{Z} \right\} \)
    \end{tasks}
}

% --- CÂU 6 ---
\questionLinesManual{0}{6}{Giá trị lớn nhất và nhỏ nhất của hàm số \( y = \dfrac{5\cos 2x + 1}{2} \) lần lượt là}{
    \begin{tasks}(4)
        \task $1$ và $2$
        \task $3$ và $2$
        \task $3$ và $-2$
        \task $-3$ và $1$
    \end{tasks}
}

% --- CÂU 7 ---
\questionLinesManual{0}{7}{Số nghiệm của phương trình \( 2\cos x - 1 = 0 \) trên đoạn \( [0; 2\pi] \) là}{
    \begin{tasks}(4)
        \task $3$
        \task $1$
        \task $4$
        \task $2$
    \end{tasks}
}

% --- CÂU 8 ---
\questionLinesManual{0}{8}{Số nghiệm của phương trình \( \dfrac{\sin 3x}{\cos x + 1} = 0 \) thuộc đoạn \( [2\pi; 4\pi] \) là}{
    \begin{tasks}(4)
        \task $7$
        \task $6$
        \task $4$
        \task $5$
    \end{tasks}
}

% --- CÂU 9 ---
\questionLinesManual{0}{9}{Tính \( \sin x \) biết \( \cos 2x = -\dfrac{3}{5} \) và \( \dfrac{\pi}{4} < x < \dfrac{\pi}{2} \).}{
    \begin{tasks}(2)
        \task \( \sin x = \dfrac{2\sqrt{5}}{5} \)
        \task \( \sin x = -\dfrac{2\sqrt{2}}{5} \)
        \task \( \sin x = \dfrac{\sqrt{17}}{5} \)
        \task \( \sin x = \dfrac{4}{5} \)
    \end{tasks}
}

% --- CÂU 10 ---
\questionLinesManual{0}{10}{Phương trình \( \cos\left(x + \dfrac{\pi}{3}\right) = 1 \) có nghiệm là:}{
    \begin{tasks}(2)
        \task \( x = -\dfrac{\pi}{3} + k2\pi \ (k \in \mathbb{Z}) \)
        \task \( x = -\dfrac{\pi}{3} + k\pi \ (k \in \mathbb{Z}) \)
        \task \( x = -\dfrac{5\pi}{6} + k2\pi \ (k \in \mathbb{Z}) \)
        \task \( x = \dfrac{\pi}{6} + k2\pi \ (k \in \mathbb{Z}) \)
    \end{tasks}
}

% --- CÂU 11 ---
\questionLinesManual{0}{11}{Số nghiệm trên đoạn \([0; 3\pi]\) của phương trình \( \sin 2x = 2\cos x \) là:}{
    \begin{tasks}(4)
        \task $4$
        \task $3$
        \task $2$
        \task $1$
    \end{tasks}
}

% --- CÂU 12 ---
\questionLinesManual{0}{12}{Cho góc \( \alpha \) thỏa \( \cot \alpha = \dfrac{3}{4} \) và \( 0^\circ < \alpha < 90^\circ \). Khẳng định nào sau đây đúng?}{
    \begin{tasks}(4)
        \task \( \cos \alpha = \dfrac{4}{5} \)
        \task \( \sin \alpha = \dfrac{4}{5} \)
        \task \( \sin \alpha = -\dfrac{4}{5} \)
        \task \( \cos \alpha = -\dfrac{4}{5} \)
    \end{tasks}
}

% --- CÂU 13 ---
\questionLinesManual{0}{13}{Tập xác định của hàm số \( y = \tan 3x \):}{
    \begin{tasks}(2)
        \task \( D = \mathbb{R} \setminus \left\{ \dfrac{\pi}{6} + k\dfrac{\pi}{3}, \, k \in \mathbb{Z} \right\} \)
        \task \( D = \mathbb{R} \setminus \left\{ \dfrac{\pi}{4} + k\pi, \, k \in \mathbb{Z} \right\} \)
        \task \( D = \mathbb{R} \setminus \left\{ \dfrac{\pi}{2} + k\pi, \, k \in \mathbb{Z} \right\} \)
        \task \( D = \mathbb{R} \setminus \left\{ k\dfrac{\pi}{2}, \, k \in \mathbb{Z} \right\} \)
    \end{tasks}
}

% --- CÂU 14 ---
\questionLinesManual{0}{14}{Giá trị lớn nhất và giá trị nhỏ nhất của hàm số \( y = 3\sin 2x - 5 \) lần lượt là:}{
    \begin{tasks}(4)
        \task \( 3; -5 \)
        \task \( -2; -8 \)
        \task \( 2; -5 \)
        \task \( 8; 2 \)
    \end{tasks}
}

% --- CÂU 15 ---
\questionLinesManual{0}{15}{Họ nghiệm của phương trình \( \sin\left(x + \dfrac{\pi}{7}\right) = \dfrac{\sqrt{3}}{2} \) là}{
    \begin{tasks}(2)
        \task 
        \(\left\{
        \begin{array}{l}
        x = \dfrac{4\pi}{21} + k2\pi \\[0.8em]
        x = \dfrac{11\pi}{21} + k2\pi 
        \end{array}
        \right. \quad (k \in \mathbb{Z})\)
        \task 
        \(\left\{
        \begin{array}{l}
        x = \dfrac{\pi}{3} + k2\pi \\[0.8em]
        x = \dfrac{2\pi}{3} + k2\pi 
        \end{array}
        \right. \quad (k \in \mathbb{Z})\)
        \task 
        \(\left\{
        \begin{array}{l}
        x = \dfrac{\pi}{3} + k2\pi \\[0.8em]
        x = -\dfrac{\pi}{3} + k2\pi 
        \end{array}
        \right. \quad (k \in \mathbb{Z})\)
        \task 
        \(\left\{
        \begin{array}{l}
        x = \dfrac{4\pi}{21} + k2\pi \\[0.8em]
        x = -\dfrac{11\pi}{21} + k2\pi 
        \end{array}
        \right. \quad (k \in \mathbb{Z})\)
    \end{tasks}
}

% --- CÂU 16 ---
\questionLinesManual{0}{16}{Tìm tập xác định \( D \) của hàm số \( y = \dfrac{3\sin x + 1}{1 - \cos 2x} \).}{
    \begin{tasks}(2)
        \task \( D = \mathbb{R} \setminus \left\{ \dfrac{\pi}{2} + k\pi \mid k \in \mathbb{Z} \right\} \)
        \task \( D = \mathbb{R} \setminus \left\{ \dfrac{\pi}{2} + k2\pi \mid k \in \mathbb{Z} \right\} \)
        \task \( D = \mathbb{R} \setminus \left\{ k2\pi \mid k \in \mathbb{Z} \right\} \)
        \task \( D = \mathbb{R} \setminus \left\{ k\pi \mid k \in \mathbb{Z} \right\} \)
    \end{tasks}
}

% --- CÂU 17 ---
\questionLinesManual{0}{17}{Giá trị lớn nhất của hàm số \( y = 4\sin x - 3 \) là}{
    \begin{tasks}(4)
        \task $-7$
        \task $-3$
        \task $1$
        \task $3$
    \end{tasks}
}

% --- CÂU 18 ---
\questionLinesManual{0}{18}{Nghiệm của phương trình \( \cos x = \dfrac{1}{2} \) là}{
    \begin{tasks}(2)
        \task \( x = \pm \dfrac{\pi}{3} + k2\pi \ (k \in \mathbb{Z}) \)
        \task \( x = \pm \dfrac{\pi}{6} + k2\pi \ (k \in \mathbb{Z}) \)
        \task \( x = \pm \dfrac{\pi}{4} + k2\pi \ (k \in \mathbb{Z}) \)
        \task \( x = \pm \dfrac{\pi}{2} + k2\pi \ (k \in \mathbb{Z}) \)
    \end{tasks}
}

% --- CÂU 19 ---
\questionLinesManual{0}{19}{Nghiệm của phương trình \( \cos x = \cos \dfrac{\pi}{12} \) là}{
    \begin{tasks}(2)
        \task 
        \(\left\{
        \begin{array}{l}
        x = \dfrac{\pi}{12} + k2\pi \\[0.8em]
        x = \dfrac{11\pi}{12} + l2\pi 
        \end{array}
        \right. \quad (k, l \in \mathbb{Z})\)
        \task 
        \(\left\{
        \begin{array}{l}
        x = \dfrac{\pi}{12} + k2\pi \\[0.8em]
        x = -\dfrac{\pi}{12} + l2\pi 
        \end{array}
        \right. \quad (k, l \in \mathbb{Z})\)
        \task 
        \(x = \dfrac{\pi}{12} + k2\pi \quad (k \in \mathbb{Z})\)
        \task 
        \(x = \dfrac{11\pi}{12} + k2\pi \quad (k \in \mathbb{Z})\)
    \end{tasks}
}

% --- CÂU 20 ---
\questionLinesManual{0}{20}{Biết \( \sin a = \dfrac{1}{3} \). Tính \( \cos 2a \) ?}{
    \begin{tasks}(4)
        \task \( \dfrac{7}{9} \)
        \task \( \dfrac{8}{9} \)
        \task \( \dfrac{2}{3} \)
        \task \( \dfrac{1}{3} \)
    \end{tasks}
}

% --- CÂU 21 ---
\questionLinesManual{0}{21}{Góc có số đo \( \dfrac{\pi}{24} \) đổi sang độ bằng}{
    \begin{tasks}(4)
        \task \( 7^\circ \)
        \task \( 7^\circ 30' \)
        \task \( 8^\circ \)
        \task \( 8^\circ 30' \)
    \end{tasks}
}

% --- CÂU 22 ---
\questionLinesManual{0}{22}{Nghiệm của phương trình \( 2\sin\left(4x - \dfrac{\pi}{3}\right) - 1 = 0 \) là:}{
    \begin{tasks}(2)
        \task \( x = \pi + k2\pi; \, x = k\dfrac{\pi}{2} \)
        \task \( x = \dfrac{\pi}{8} + k\dfrac{\pi}{2}; \, x = \dfrac{7\pi}{24} + k\dfrac{\pi}{2} \)
        \task \( x = k2\pi; \, x = \dfrac{\pi}{2} + k2\pi \)
        \task \( x = k\pi; \, x = \pi + k2\pi \)
    \end{tasks}
}

% --- CÂU 23 ---
\questionLinesManual{0}{23}{Phương trình \( \sin x = \dfrac{\sqrt{3}}{2} \) có nghiệm là:}{
    \begin{tasks}(2)
        \task \( x = \pm \dfrac{\pi}{3} + k2\pi \)
        \task \( x = \dfrac{\pi}{3} + k\pi \)
        \task 
        \(\left\{
        \begin{array}{l}
        x = \dfrac{\pi}{6} + k\pi \\[0.8em]
        x = \dfrac{5\pi}{6} + k\pi 
        \end{array}
        \right.\)
        \task 
        \(\left\{
        \begin{array}{l}
        x = \dfrac{\pi}{3} + k2\pi \\[0.8em]
        x = \dfrac{2\pi}{3} + k2\pi 
        \end{array}
        \right.\)
    \end{tasks}
}
% --- CÂU 24 ---
\questionLinesManual{0}{24}{Tìm tất cả các nghiệm của phương trình \( \sin\left(x + \dfrac{\pi}{6}\right) = 1 \).}{
    \begin{tasks}(2)
        \task \( x = \dfrac{\pi}{3} + k\pi \, (k \in \mathbb{Z}) \)
        \task \( x = -\dfrac{\pi}{6} + k2\pi \, (k \in \mathbb{Z}) \)
        \task \( x = \dfrac{\pi}{3} + k2\pi \, (k \in \mathbb{Z}) \)
        \task \( x = \dfrac{5\pi}{6} + k2\pi \, (k \in \mathbb{Z}) \)
    \end{tasks}
}

% --- CÂU 25 ---
\questionLinesManual{0}{25}{Phương trình \( 2\sin x - 1 = 0 \) có tập nghiệm là:}{
    \begin{tasks}(2)
        \task \( S = \left\{ \dfrac{\pi}{6} + k2\pi; \, \dfrac{5\pi}{6} + k2\pi, \, k \in \mathbb{Z} \right\} \)
        \task \( S = \left\{ \dfrac{\pi}{3} + k2\pi; \, -\dfrac{2\pi}{3} + k2\pi, \, k \in \mathbb{Z} \right\} \)
        \task \( S = \left\{ \dfrac{\pi}{6} + k2\pi; \, -\dfrac{\pi}{6} + k2\pi, \, k \in \mathbb{Z} \right\} \)
        \task \( S = \left\{ \dfrac{1}{2} + k2\pi, \, k \in \mathbb{Z} \right\} \)
    \end{tasks}
}

% --- CÂU 26 ---
\questionLinesManual{0}{26}{Phương trình \( 2\cos x - \sqrt{2} = 0 \) có tất cả các nghiệm là}{
    \begin{tasks}(2)
        \task 
        \(\left\{
        \begin{array}{l}
        x = \dfrac{3\pi}{4} + k2\pi \\[0.8em]
        x = -\dfrac{3\pi}{4} + k2\pi 
        \end{array}
        \right. \quad (k \in \mathbb{Z})\)
        \task 
        \(\left\{
        \begin{array}{l}
        x = \dfrac{\pi}{4} + k2\pi \\[0.8em]
        x = -\dfrac{\pi}{4} + k2\pi 
        \end{array}
        \right. \quad (k \in \mathbb{Z})\)
        \task 
        \(\left\{
        \begin{array}{l}
        x = \dfrac{\pi}{4} + k2\pi \\[0.8em]
        x = \dfrac{3\pi}{4} + k2\pi 
        \end{array}
        \right. \quad (k \in \mathbb{Z})\)
        \task 
        \(\left\{
        \begin{array}{l}
        x = \dfrac{7\pi}{4} + k2\pi \\[0.8em]
        x = -\dfrac{7\pi}{4} + k2\pi 
        \end{array}
        \right. \quad (k \in \mathbb{Z})\)
    \end{tasks}
}

% --- CÂU 27 ---
\questionLinesManual{0}{27}{Giải phương trình \( 2\cos x - 1 = 0 \)}{
    \begin{tasks}(2)
        \task \( x = \pm \dfrac{\pi}{3} + k\pi, \, k \in \mathbb{Z} \)
        \task 
        \(\left\{
        \begin{array}{l}
        x = \dfrac{\pi}{3} + k2\pi \\[0.8em]
        x = \dfrac{2\pi}{3} + k2\pi 
        \end{array}
        \right. \quad (k \in \mathbb{Z})\)
        \task \( x = \pm \dfrac{\pi}{3} + k2\pi, \, k \in \mathbb{Z} \)
        \task 
        \(\left\{
        \begin{array}{l}
        x = \dfrac{\pi}{3} + k\pi \\[0.8em]
        x = \dfrac{2\pi}{3} + k\pi 
        \end{array}
        \right. \quad (k \in \mathbb{Z})\)
    \end{tasks}
}

% --- CÂU 28 ---
\questionLinesManual{0}{28}{Nghiệm của phương trình \( \cos x = -1 \) là:}{
    \begin{tasks}(2)
        \task \( x = \dfrac{\pi}{2} + k\pi, \, k \in \mathbb{Z} \)
        \task \( x = k2\pi, \, k \in \mathbb{Z} \)
        \task \( x = \pi + k2\pi, \, k \in \mathbb{Z} \)
        \task \( x = k\pi, \, k \in \mathbb{Z} \)
    \end{tasks}
}

% --- CÂU 29 ---
\questionLinesManual{0}{29}{Phương trình lượng giác: \( 2\cos x + \sqrt{2} = 0 \) có nghiệm là}{
    \begin{tasks}(2)
        \task 
        \(\left\{
        \begin{array}{l}
        x = \dfrac{\pi}{4} + k2\pi \\[0.8em]
        x = -\dfrac{\pi}{4} + k2\pi 
        \end{array}
        \right.\)
        \task 
        \(\left\{
        \begin{array}{l}
        x = \dfrac{3\pi}{4} + k2\pi \\[0.8em]
        x = -\dfrac{3\pi}{4} + k2\pi 
        \end{array}
        \right.\)
        \task 
        \(\left\{
        \begin{array}{l}
        x = \dfrac{\pi}{4} + k2\pi \\[0.8em]
        x = \dfrac{3\pi}{4} + k2\pi 
        \end{array}
        \right.\)
        \task 
        \(\left\{
        \begin{array}{l}
        x = \dfrac{7\pi}{4} + k2\pi \\[0.8em]
        x = -\dfrac{7\pi}{4} + k2\pi 
        \end{array}
        \right.\)
    \end{tasks}
}

% --- CÂU 30 ---
\questionLinesManual{0}{30}{Phương trình lượng giác \( 3\cot x - \sqrt{3} = 0 \) có nghiệm là:}{
    \begin{tasks}(2)
        \task \( x = \dfrac{\pi}{3} + k2\pi \)
        \task Vô nghiệm
        \task \( x = \dfrac{\pi}{6} + k\pi \)
        \task \( x = \dfrac{\pi}{3} + k\pi \)
    \end{tasks}
}

\end{document}